\chapter{Dictionaries in Python: Key-Value Data Storage}

\section{Introduction: Real-Life Hook}

Students often deal with pairs of information:

\begin{itemize}
    \item Student name $\rightarrow$ roll number  
    \item Phonebook: name $\rightarrow$ phone number  
    \item Grocery items $\rightarrow$ price  
\end{itemize}

Python’s \textbf{dictionary} stores data as \textbf{key-value pairs}, making it easier to look up information quickly.

\section{Concept: What is a Dictionary?}

\begin{definition}{Dictionary}
A \textbf{dictionary} is a collection of key-value pairs in Python. Each key is unique, and it is used to access its corresponding value.
\end{definition}

\subsection{Creating a Dictionary}

\begin{verbatim}
student = {
    "name": "Amina",
    "roll": 12,
    "marks": 85
}
\end{verbatim}

- Curly braces \{\} are used  
- Each entry is a \textbf{key: value} pair  
- Keys are unique; values can be repeated

\subsection{Accessing Values}

\begin{verbatim}
print(student["name"])  # Amina
print(student["marks"]) # 85
\end{verbatim}

\subsection{Updating Values}

\begin{verbatim}
student["marks"] = 90
print(student)  # {'name': 'Amina', 'roll': 12, 'marks': 90}
\end{verbatim}

\subsection{Adding New Key-Value Pairs}

\begin{verbatim}
student["age"] = 15
print(student)
\end{verbatim}

\subsection{Removing Key-Value Pairs}

\begin{verbatim}
del student["roll"]
print(student)
\end{verbatim}

\section{Looping Through a Dictionary}

\begin{verbatim}
student = {"name": "Amina", "marks": 90, "age": 15}

# Loop through keys
for key in student:
    print(key, ":", student[key])

# Loop through key-value pairs
for key, value in student.items():
    print(key, "->", value)
\end{verbatim}

Output:

\begin{verbatim}
name : Amina
marks : 90
age : 15
name -> Amina
marks -> 90
age -> 15
\end{verbatim}

\section{Nested Dictionaries}

\begin{verbatim}
students = {
    "s1": {"name": "Amina", "marks": 90},
    "s2": {"name": "Rafi", "marks": 85}
}

print(students["s1"]["name"])  # Amina
\end{verbatim}

- Dictionaries can contain other dictionaries for complex data

\section{Mini-Project: Phonebook Program}

Create a program to store phone numbers and look them up:

\begin{verbatim}
phonebook = {
    "Amina": "017XXXXXXXX",
    "Rafi": "018XXXXXXXX",
    "Sadia": "019XXXXXXXX"
}

name = input("Enter a name to search: ")
if name in phonebook:
    print("Phone number:", phonebook[name])
else:
    print("Name not found")
\end{verbatim}

\section{Step-by-Step Diagram: Dictionary Lookup Flow}

\begin{center}
\begin{tikzpicture}[node distance=1.5cm]
\node (start) [draw, ellipse, fill=green!20] {Start Program};
\node (input) [draw, rectangle, fill=yellow!20, below of=start] {Get Key from User};
\node (check) [draw, diamond, fill=orange!20, below of=input] {Key Exists?};
\node (yes) [draw, rectangle, fill=yellow!20, below left of=check, xshift=-2cm] {Print Value};
\node (no) [draw, rectangle, fill=yellow!20, below right of=check, xshift=2cm] {Print "Not Found"};
\node (end) [draw, ellipse, fill=green!20, below of=check, yshift=-3cm] {End Program};
\draw[->] (start) -- (input);
\draw[->] (input) -- (check);
\draw[->] (check.west) -| node[near start, above] {Yes} (yes.north);
\draw[->] (check.east) -| node[near start, above] {No} (no.north);
\draw[->] (yes.south) |- (0,-7) -| (end.north);
\draw[->] (no.south) |- (0,-7) -| (end.north);
\end{tikzpicture}
\end{center}

\section{Key Concepts}

\begin{itemize}
    \item Dictionaries store \textbf{key-value pairs} for fast lookup  
    \item Keys are unique; values can be anything  
    \item Dictionaries can be \textbf{nested} for structured data  
    \item Loops can iterate over keys, values, or items  
\end{itemize}

\section{End-of-Chapter Exercises}

\begin{exercise}{Very Short Questions}
\begin{enumerate}
    \item What is a dictionary in Python?  
    \item How do you access a value using a key?  
    \item Can keys in a dictionary be repeated?  
\end{enumerate}
\end{exercise}

\begin{exercise}{Short Answer Questions}
\begin{enumerate}
    \item Create a dictionary to store 3 students’ marks and print each student’s marks.  
    \item How do you delete a key-value pair from a dictionary?  
    \item Explain the difference between a list and a dictionary.  
\end{enumerate}
\end{exercise}

\begin{exercise}{Activity-Based Questions}
\begin{enumerate}
    \item Create a dictionary of 5 countries and their capitals. Allow the user to look up a capital by country name.  
    \item Make a nested dictionary to store 2 students, each with name, age, and marks. Print their information using loops.  
    \item Create a dictionary of items and prices. Calculate the total price of all items using a loop.  
\end{enumerate}
\end{exercise}

\section{Teacher Notes}

\marginnote{
\textbf{Learning Objectives:}  
\begin{itemize}
    \item Students understand dictionary creation, key-value access, and updates  
    \item Students can use loops to iterate through dictionaries  
    \item Students can create nested dictionaries for complex data  
\end{itemize}

\textbf{Teaching Tips:}  
\begin{itemize}
    \item Use real-life examples like phonebooks, students’ marks, grocery prices  
    \item Demonstrate dictionary lookup and updates live in class  
\end{itemize}

\textbf{Assessment:}  
\begin{itemize}
    \item Check correctness of dictionary operations  
    \item Observe understanding of keys, values, loops, and nested dictionaries  
\end{itemize}
}